% ============================================================================
\chapter{Frameworks PHP~: Laravel, Symfony, et les autres}
\label{ch-php-frameworks}
% ============================================================================

\begin{objectif}
\begin{itemize}
  \item Situer les cadres PHP en 2026, et choisir sans dogme.
  \item Reconna\^itre leurs concepts communs~: routage, DI, ORM, migrations, vues, files, tests.
  \item \'Ecrire, dans chacun, le CRUD qu'on vient de faire \og ~\`a la main~ \fg{}, et mesurer le gain.
  \item Comprendre les templates et l'\'ecosyst\`eme (bundles, packages, Forge, Octane...).
  \item Savoir quand s'en passer, et comment migrer une application legacy vers l'un d'eux.
\end{itemize}
\end{objectif}

\begin{note}
Ce chapitre n'est pas un cours de Laravel ni de Symfony~: ces cadres ont chacun une
documentation vaste et un livre entier leur serait consacr\'e. Il te donne la \textbf{carte}~:
tu sauras o\`u chercher, et surtout, tout ce que tu as \'ecrit \og ~\`a la main~ \fg{} dans les chapitres
pr\'ec\'edents porte un nom dans ces cadres. C'est cet aller-retour~: comprendre le m\'ecanisme,
puis utiliser l'outil.
\end{note}

\section{Le paysage}

\begin{tabz}{lXll}
  \headrow \thd{Cadre} \& \thd{Ce qui le d\'efinit} \& \thd{Version} \& \thd{Lien} \\
  \midrule
  \link{Laravel}{https://laravel.com/} \& la productivit\'e~: Eloquent, Artisan, \'ecosyst\`eme \'enorme (Nova, Forge, Livewire, Inertia), tr\`es bonne documentation, le plus d'offres \og ~web agency~ \fg{}. \& 12 \& \url{https://laravel.com/docs} \\
  \link{Symfony}{https://symfony.com/} \& les composants r\'eutilisables + un cadre~: rigueur, DI explicite, long support LTS, tr\`es pr\'esent en France/Europe et dans le secteur public. \& 7.x (6.4 LTS) \& \url{https://symfony.com/doc/current/index.html} \\
  \link{Slim}{https://www.slimframework.com/} \& un routeur + PSR-7/15, rien d'autre~: id\'eal pour une API l\'eg\`ere. \& 4.x \& \url{https://www.slimframework.com/docs/} \\
  \link{Laminas / Mezzio}{https://mezzio.dev/} \& PSR pur, modulaire, tr\`es typ\'e entreprise. \& --- \& \url{https://docs.laminas.dev/} \\
  \link{CodeIgniter~4}{https://codeigniter.com/} \& l\'eger, simple, h\'eritage du mutualis\'e. \& 4.x \& \url{https://codeigniter.com/user_guide/} \\
  \link{Yii~2}{https://www.yiiframework.com/} \& Gii (g\'en\'eration de CRUD), tr\`es rapide, populaire en Asie. \& 2.0 \& \url{https://www.yiiframework.com/doc/} \\
  \link{Phalcon}{https://phalcon.io/} \& extension C~: performances brutes \'elev\'ees. \& 5.x \& \url{https://docs.phalcon.io/} \\
  \link{CakePHP}{https://book.cakephp.org/} \& \og ~convention over configuration~ \fg{}, le doyen. \& 5.x \& \\
  \link{WordPress}{https://wordpress.org/develop/} \& pas un cadre MVC, mais 40 \% du web~: l'API Plugin/Theme vaut le d\'etour pro. \& 6.x \& \url{https://developer.wordpress.org/} \\
\end{tabz}

\section{Le dictionnaire commun~: le m\^eme objet, trois noms}

\begin{tabz}{lXXX}
  \headrow \thd{Concept} \& \thd{En PHP nu (tes chapitres)} \& \thd{Laravel} \& \thd{Symfony} \\
  \midrule
  Point d'entr\'ee unique \& \verb|public/index.php| \& \verb|public/index.php| + \verb|bootstrap/app.php| \& \verb|public/index.php| + \verb|Kernel| \\
  Routage \& ta table \verb|\$routes| \& \verb|routes/web.php|, \verb|Route::get()| \& attributs \verb|\#[Route]| sur un contr\^oleur \\
  Contr\^oleur \& ta fonction \& \verb|app/Http/Controllers| \& \verb|src/Controller| \\
  Injection de d\'ependances \& ton conteneur maison \& le conteneur/service container auto-wired \& \verb|services.yaml| + autoconfigure \\
  ORM \& ton d\'ep\^ot PDO \& Eloquent (ActiveRecord) \& Doctrine ORM (Data Mapper) \\
  Migrations \& Phinx \& \verb|php artisan make:migration| \& \verb|bin/console doctrine:migrations:migrate| \\
  Validation \& ton \verb|Valideur| \& Form Requests + \verb|validate()| \& symfony/validator + attributs \\
  Vues \& \verb|require| + \'echappement \& Blade \& Twig \\
  Journal \& \verb|error\_log| \& \verb|Log::channel()| \& Monolog (c'est un composant Symfony) \\
  Files d'attente \& Redis + cron \& \verb|dispatch()| + Horizon \& Messenger + transport \\
  T\^aches planifi\'ees \& crontab \& scheduler \& Messenger scheduled tasks \\
  Tests \& PHPUnit \& \verb|php artisan test| (Pest par d\'efaut) \& PHPUnit + \verb|WebTestCase|/Panther \\
\end{tabz}

\subsection{Le m\^eme CRUD dans les trois styles}

\begin{codephp}
<?php
// ---------- LARAVEL : routes/api.php
use App\Http\Controllers\ArticleController;
Route::apiResource('articles', ArticleController::class)->parameters(['articles' => 'slug']);

// ---------- LARAVEL : ArticleController.php
namespace App\Http\Controllers;

use App\Http\Requests\StoreArticle;
use App\Models\Article;

class ArticleController extends Controller
{
    public function index(): \Illuminate\Http\Resources\Json\AnonymousResourceCollection
    {
        return Article::with('auteur:id,nom,email')          // evite le N+1
            ->where('statut', 'publie')
            ->latest('publie_le')
            ->paginate(20)
            ->through(fn (Article $a) => ArticleResource::make($a));
    }

    public function store(StoreArticle $request): \Symfony\Component\HttpFoundation\Response
    {
        $article = Article::create($request->validated() + ['auteur_id' => $request->user()->id]);
        return response()->json(ArticleResource::make($article), 201)
            ->header('Location', route('articles.show', $article));
    }
}

// ---------- LARAVEL : StoreArticle.php (la validation devient un objet reutilisable)
class StoreArticle extends \Illuminate\Foundation\Http\FormRequest
{
    public function rules(): array
    {
        return [
            'titre' => ['required', 'string', 'min:3', 'max:160'],
            'corps' => ['required', 'string', 'min:20', 'max:100000'],
            'tags' => ['array', 'max:12'],
            'tags.*' => ['string', 'max:30', 'distinct'],
        ];
    }
}

// ---------- LARAVEL : Models/Article.php (ActiveRecord : le modele porte la persistance)
class Article extends \Illuminate\Database\Eloquent\Model
{
    protected $fillable = ['titre', 'corps', 'slug', 'statut', 'publie_le', 'auteur_id'];
    protected $casts = ['publie_le' => 'datetime', 'tags' => 'array'];

    public function auteur(): \Illuminate\Database\Eloquent\Relations\BelongsTo
    { return $this->belongsTo(Auteur::class); }
    public function commentaires(): \Illuminate\Database\Eloquent\Relations\HasMany
    { return $this->hasMany(Commentaire::class)->where('approuve', true); }

    public function scopePublies($q) { return $q->where('statut', 'publie'); }
}
\end{codephp}

\begin{codephp}
<?php
// ---------- SYMFONY : src/Controller/ArticleController.php
namespace App\Controller;

use App\Entity\Article;
use App\Repository\ArticleRepository;
use App\Form\ArticleType;
use Doctrine\ORM\EntityManagerInterface;
use Symfony\Bundle\FrameworkBundle\Controller\AbstractController;
use Symfony\Component\HttpFoundation\{Request, Response, JsonResponse};
use Symfony\Component\Routing\Attribute\{Route, Method};
use Symfony\Component\Validator\Constraints as Assert;

final class ArticleController extends AbstractController
{
    public function __construct(private readonly ArticleRepository $articles) {}

    #[Route('/api/articles', name: 'api_articles', methods: ['GET'])]
    public function index(Request $request): JsonResponse
    {
        $page = max(1, (int) $request->query->getInt('page', 1));
        $items = $this->articles->parPage($page, min(100, max(1, $request->query->getInt('limite', 20))));
        return $this->json([
            'items' => $items,
            'page' => $page,
            'total' => $this->articles->compterPublies(),
        ], headers: ['cache-control' => 'public, max-age=30']);
    }

    #[Route('/articles/new', name: 'article_new', methods: ['GET', 'POST'])]
    public function new(Request $request, EntityManagerInterface $em): Response
    {
        $article = new Article();
        $form = $this->createForm(ArticleType::class, $article);
        $form->handleRequest($request);
        if ($form->isSubmitted() && $form->isValid()) {
            $article->publier(new \DateTimeImmutable());
            $em->persist($article);
            $em->flush();                       // une seule ecriture, la transaction est implicite
            $this->addFlash('succes', 'Article cree.');
            return $this->redirectToRoute('article_voir', ['slug' => $article->getSlug()]);
        }
        return $this->render('article/new.html.twig', ['form' => $form]);
    }
}

// ---------- SYMFONY : src/Entity/Article.php (Data Mapper : l'entite est pure)
#[\Doctrine\ORM\Mapping\Entity]
#[\Doctrine\ORM\Mapping\Table(name: 'articles')]
#[\Doctrine\ORM\Mapping\HasLifecycleCallbacks]
class Article
{
    #[\Doctrine\ORM\Mapping\Id, \Doctrine\ORM\Mapping\GeneratedValue, \Doctrine\ORM\Mapping\Column]
    private ?int $id = null;

    #[Assert\Length(min: 3, max: 160)]
    #[\Doctrine\ORM\Mapping\Column(length: 191, unique: true)]
    private string $slug = '';

    /** @phpstan-ignore-next-line */
    #[\Doctrine\ORM\Mapping\ManyToOne(targetEntity: Utilisateur::class)]
    #[\Doctrine\ORM\Mapping\JoinColumn(nullable: false)]
    private ?Utilisateur $auteur = null;

    public function publier(\DateTimeImmutable $quand): void
    {
        $this->publieLe = $quand;
        $this->statut = Statut::Publie;
    }
    /* getters/setters : php bin/console make:entity les genere */
}
\end{codephp}

\begin{codephp}
<?php
// ---------- Slim 4 : le "presque rien" pour une API
use Psr\Http\Message\ServerRequestInterface as Requete;
use Psr\Http\Message\ResponseInterface as Reponse;
use Slim\Factory\AppFactory;

$app = AppFactory::create();
$app->addBodyParsingMiddleware();
$app->addRoutingMiddleware();
$app->addErrorMiddleware(true, true, true);          // affiche les details en dev

$app->get('/api/v1/articles', function (Requete $r, Reponse $res, array $args) use ($depot) {
    $q = $r->getQueryParams();
    $donnees = $depot->page((int) ($q['page'] ?? 1), min(100, (int) ($q['limite'] ?? 20)));
    $res->getBody()->write(json_encode($donnees, JSON_THROW_ON_ERROR | JSON_UNESCAPED_UNICODE));
    return $res->withHeader('content-type', 'application/json');
});
$app->run();
\end{codephp}

\section{Les templates~: Blade et Twig}

\begin{codeplain}
{{-- Blade (Laravel) --}}
@extends('layouts.app')
@section('title', $article->titre)

@section('content')
    <article>
        <h1>{{ $article->titre }}</h1>          {{-- echappe par defaut ; {!! !!} = brut (dangereux) --}}
        <p>{{ Str::limit($article->extrait, 200) }}</p>
        @can('update', $article)
            <a href="{{ route('articles.edit', $article) }}">Modifier</a>
        @endcan
        @foreach ($article->tags as $tag)
            <x-ui.badge>{{ $tag->nom }}</x-ui.badge>       {{-- un composant Blade --}}
        @endforeach
        @empty($article->commentaires)
            <p>Aucun commentaire.</p>
        @endempty
    </article>
@endsection
\end{codeplain}

\begin{codeplain}
{# Twig (Symfony) #}
{% extends 'base.html.twig' %}
{% block title %}{{ article.titre }}{% endblock %}
{% block content %}
  <article>
    <h1>{{ article.titre }}</h1>                          {# echappe par defaut #}
    <p>{{ article.extrait|u.truncate(200, '...') }}</p>
    {% is_granted('update', article) %}
      <a href="{{ path('article_edit', { slug: article.slug }) }}">Modifier</a>
    {% endis_granted %}
    {% for tag in article.tags %}
      {% include 'components/_badge.html.twig' with { label: tag.nom } %}
    {% else %}
      <p>Aucun tag.</p>
    {% endfor %}
    {{ form_start(form) }}{{ form_end(form) }}
  </article>
{% endblock %}
\end{codeplain}

\section{Les outils qui accompagnent chaque cadre}

\begin{tabz}{XX}
  \headrow \thd{Laravel} \& \thd{Symfony} \\
  \midrule
  \verb|artisan| (g\'en\'erateurs, migrations, tinker) \& \verb|bin/console| (+ debug:router, debug:container) \\
  \link{Eloquent}{https://laravel.com/docs/eloquent} \& \link{Doctrine}{https://www.doctrine-project.org/projects/orm.html} \\
  \link{Livewire}{https://livewire.laravel.com/} / \link{Inertia}{https://inertiajs.com/} (interactive sans SPA) \& \link{Stimulus + Turbo (Hotwire)}{https://symfony.com/doc/current/frontend/htmx.html} \\
  \link{Forge}{https://forge.laravel.com/} / \link{Envoyer}{https://envoyer.io/} (d\'eploiement) \& \link{Symfony~Cloud}{https://symfony.com/cloud/} (PaaS officiel) \\
  \link{Octane}{https://laravel.com/docs/octane} (FrankenPHP/Swoole, ~5x d\'ebit) \& \link{FrankenPHP}{https://frankenphp.dev/} (app worker, idem, et c'est un serveur web \`a lui seul) \\
  Nova (back-office), Horizon (files), Reverb (WebSocket) \& Blackfire, Messenger, UX, Flex, MakerBundle \\
  \link{Name~of~Zen~docs}{https://laravel-news.com/} \& \url{https://symfony.com/blog/} + la newsletter \\
\end{tabz}

\begin{astuce}
\textbf{FrankenPHP} m\'erite ton attention~: un serveur applicatif PHP (fork de Caddy) qui garde le
processus vivant~: d\'emarrage d'app quasi instantan\'e, mont\'ee en d\'ebit, mont\'ees de charge \og ~\`a la
Node~ \fg{}, et int\'egration avec Mercure/SSE. C'est la nouveaut\'e d'infrastructure PHP la plus
int\'eressante depuis longtemps, et elle marche avec Laravel (Octane), Symfony et Slim.
\end{astuce}

\section{Choisir --- et migrer}

\begin{enumerate}
  \item \textbf{Tu veux livrer vite, seul, avec peu de lignes~?} Laravel. La communaut\'e est tr\`es
        active, presque tout existe d\'ej\`a en paquet.
  \item \textbf{Tu vises une entreprise, un march\'e fran\c{c}ais/europ\'een, un projet \`a longue dur\'ee de
        vie, ou tu aimes le typage fort~?} Symfony. Ses composants sont utilisables sans le cadre.
  \item \textbf{API minuscule, ou tu veux garder la main~?} Slim, ou ton noyau du chapitre~7.
  \item \textbf{Site \'editorial, clients non techniques~?} WordPress (avec Roots Bedrock + Acorn si tu veux du moderne). C'est un m\'etier \`a part~; ne le m\'eprise pas, mais ne l'appelle pas un cadre.
  \item \textbf{Portfolio et apprentissage~?} Fais le m\^eme projet en PHP nu, puis en Laravel~: tu
        sauras expliquer \emph{pourquoi} le cadre est utile. C'est exactement ce qui distingue un
        d\'eveloppeur d'un utilisateur de tutoriel.
\end{enumerate}

\etape{Migration d'un projet legacy~: la strat\'egie r\'ealiste}
\begin{enumerate}
  \item Ne r\'e\'ecris rien d'abord~: ajoute \verb|declare(strict\_types=1)|, Composer, PSR-4,
        PHPStan, et les tests sur le comportement actuel (tests \og ~de caract\'erisation~ \fg{}).
  \item Extrais le routeur et la liaison vers les contr\^oleurs (une table, une fonction par page).
  \item Introduis un conteneur de service~; remplace les \verb|new| globaux et les
        \verb|include| par de l'injection~; enrobe l'acc\`es \`a la base dans un d\'ep\^ot.
  \item Remplace les \verb|echo| \'epars par un moteur de template (Twig~: il s'installe seul).
  \item Garde l'application \og ~enrob\'ee~ \fg{} dans un cadre (Symfony/Laravel peuvent servir un script
        legacy via une route)~: les pages migrent une par une, pendant des mois, sans gros plan.
\end{enumerate}

\begin{exercice}
\begin{enumerate}
  \item Installe Laravel (\cmd{composer create-project laravel/laravel blog}) et Symfony
        (\cmd{composer create-project symfony/skeleton:"7.2.\*" blog}), fais tourner les deux,
        et trouve la route d'accueil de chacun dans le code~: tu dois pouvoir dire o\`u le
        fichier\ est et pourquoi.
  \item Dans Laravel~: mod\`ele \verb|Article| + migration + factory + seeder, contr\^oleur resource,
        validation par FormRequest, ressource JSON, paginate, autorisation avec Policies~;
        10 tests (feature tests avec SQLite en m\'emoire).
  \item Dans Symfony~: m\^eme chose avec MakerBundle (\verb|make:entity|, \verb|make:controller|),
        Doctrine fixtures, formulaire, validation par attributs, un test fonctionnel avec
        \verb|WebTestCase|.
  \item Compare la production~: nombre de fichiers, lignes \'ecrites, temps pass\'e, courbe
        d'apprentissage~; r\'edige une d\'ecision argument\'ee dans \verb|ANALYSE.md| (c'est un livrable
        que les recruteurs lisent).
  \item Ajoute dans le projet Laravel une file d'attente~: \`a la publication, envoyer un email en
        arri\`ere-plan avec \verb|ShouldQueue|, un job retry et un \'echec visible dans Horizon.
  \item D\'eploie le projet Symfony ou Laravel sur un VPS avec Deployer ou Ploi~: z\'ero
        \verb|git pull| manuel, migrations automatiques, OPcache invalid\'e, HTTPS, et rollback.
\end{enumerate}
\end{exercice}

\begin{ressources}
\begin{itemize}
  \item Laravel --- docs, Laracasts, News --- \url{https://laravel.com/docs/12.x}, \url{https://laracasts.com/}, \url{https://laravel-news.com/}
  \item Symfony --- les guides (excellents, en fran\c{c}ais partiel) --- \url{https://symfony.com/doc/current/page_blocks.html} \\
        et la formation gratuite Open Classrooms --- \url{https://openclassrooms.com/fr/search?query=symfony}
  \item \link{Symfony~Casts}{https://symfonycasts.com/} (les meilleures vid\'eos de l'\'ecosyst\`eme, beaucoup sont gratuites)
  \item \link{Laravel~Daily~Learn}{https://laraveldaily.com/learn}, \link{Relay}{https://laravelrelay.app/}
  \item Comparatifs honn\^etes et r\'ecents --- \url{https://packagist.org/stats/} (parts d'usage des paquets)
\end{itemize}
\end{ressources}

\index{Laravel}\index{Symfony}\index{Blade}\index{Twig}\index{migration legacy}
